\section{Axis 3: Computational Efficiency}
\label{sec:efficiency}

The computational bottleneck of neighbourhood rough sets is the dependency
evaluation $\DEP(B \cup \{a\})$ for each candidate attribute $a$ at each
forward-greedy step.  Na\"ively, this costs $O(n^2 d)$: for each of the
$n^2/2$ object pairs, check whether they share a neighbourhood granule
under $B$.  For large $n$ (thousands of objects) or large $d$ (gene expression),
this is prohibitive.  Three families of acceleration have emerged.

%-----------------------------------------------------------------------
\subsection{Object-Pair Acceleration}
\label{sec:pair_accel}
%-----------------------------------------------------------------------

\paragraph{Maximal discernibility pairs (Dai et al., 2018).}
Rather than maintaining a dependency degree over all $n^2/2$ pairs, Dai et
al.~\cite{Dai2018TFS} observe that once a pair $(x_i, x_j)$ has been
\emph{discerned} by the current attribute subset $B$ (i.e., they fall in
different neighbourhoods), it does not need to be revisited when adding
further attributes.  The working set of pairs shrinks monotonically as $B$
grows.  This yields the concept of \emph{reduced maximal discernibility
pairs}: the minimal set of cross-class object pairs sufficient to determine
which attributes can further increase $\DEP$.  Empirically, the working set
reaches $\leq 20\%$ of all pairs on typical UCI datasets.

\paragraph{Parallel discernibility matrix (Sowkuntla \& Prasad Rao, 2021).}
The discernibility matrix representation of a fuzzy rough set---one row per
cross-class pair, one column per attribute---is sparser than the full
similarity matrix and decomposes naturally across rows.
Sowkuntla and Prasad Rao~\cite{Sowkuntla2021DARA} implement this decomposition
via MapReduce, reporting linear speedup on large datasets.  The key observation---discernibility
matrix rows are cheaper and more parallelisable than similarity matrix rows---was
identified independently but shares its core argument with Dai et al.'s
working-set reduction.

\paragraph{Discernibility hashing (Luo et al., 2025).}
Luo et al.~\cite{Luo2025hash} further compress the discernibility matrix by
hashing each row's attribute support set into a fixed-size hash space.
Two rows with the same support hash to the same bucket; only canonical
representatives per bucket are processed.  This reduces the effective matrix
size from $O(n^2)$ rows to $O(H)$ buckets (where $H$ is the hash table size),
with controllable collision probability.

%-----------------------------------------------------------------------
\subsection{Feature-Space Acceleration}
\label{sec:feat_accel}
%-----------------------------------------------------------------------

\paragraph{Grouping + representative selection.}
A complementary strategy groups attributes by mutual redundancy (clustering in
feature space) and selects one representative per group.  Kuzudisli et
al.~\cite{Kuzudisli2023review} survey this family, which encompasses FCBF
(Yu \& Liu~\cite{YuLiu2004}, 2321 citations), graph-based grouping
(Zheng et al.~\cite{Zheng2021grouping}; Wan et al.~\cite{Wan2023TFS}), and
cascaded fuzzy clustering (Chen et al.~\cite{Chen2024cascade}).
The redundancy within each group can be formalised in granular terms:
two attributes $a, b$ are \emph{granularly redundant} if
$\NB{\{a\}}(x) = \NB{\{b\}}(x)$ for all $x$; the granule geometry is then
a natural basis for the clustering step.

%-----------------------------------------------------------------------
\subsection{Distance Concentration as a Barrier}
\label{sec:concentration}
%-----------------------------------------------------------------------

All pair-pruning strategies share a vulnerability at high dimension:
\emph{concentration of measure}.  In $d$-dimensional space, the ratio
$(\max_{i,j} \rho(x_i,x_j)) / (\min_{i,j} \rho(x_i,x_j)) \to 1$ as
$d \to \infty$ for many distributions.  When all distances are nearly equal,
pruning by triangle inequality (``if $\rho(a,c) > \rho(a,b) + \rho(b,c)$
then prune'') offers no saving because no triangle inequality is violated.
Song et al.'s TRIM~\cite{TRIM2025} showed in the vector-search context that
pruning ratios collapse to near-zero at $d > 100$ without landmark-based
corrections.

The same barrier applies to rough set pair acceleration at high $d$.
The literature benchmarks primarily on $d \leq 60$ (Sonar, Ionosphere);
GBNRS reports results on up to $d = 166$ (Musk1).  Whether any of these
acceleration techniques survive at $d \sim 10^3$ (genomics) is unknown
and constitutes an open problem (Section~\ref{sec:open}).
